% ==================== Chapter 5: Testing and Evaluation ====================
\newpage\section{System Testing, Evaluation and Validation} [Estimated: 5,500 words]

This chapter presents a comprehensive testing and evaluation strategy for the Android permission auditing and compliance early-warning framework proposed in this study. The testing regime spans multiple levels, from unit-level logic correctness to end-to-end real-device scenarios, with a particular emphasis on a stochastic stress-testing methodology based on an automated violation simulator. This approach enables scientifically rigorous, repeatable validation of system robustness under complex and unpredictable real-world usage conditions. All tests follow the principles of test-driven development and continuous integration, ensuring that code changes do not introduce regressions. The sections below are organised from fine-grained to coarse-grained testing.

\subsection{Unit Testing and Baseline Logic Correctness} [Estimated: 1,000 words]

Unit testing forms the foundation of the validation pyramid, aiming to ensure that each functional module operates as expected in isolation. This study employs JUnit 5 as the primary unit testing framework, supplemented by Mockito for dependency mocking, thereby eliminating interference from external environments (such as Android system services and databases) on test outcomes. Test cases cover the following three core logic components:

\begin{enumerate}
\item \textbf{Baseline parsing and mapping}: Verify that the system can correctly load and parse predefined baseline configuration files (for the four categories: MAPS, SOCIAL, TOOLS, DEFAULT). A range of both valid and malformed JSON configuration files are prepared; the parser is checked for accurate extraction of package names, baseline frequencies, and standard deviations, and for throwing clear exceptions upon format errors. Additionally, the mapping logic—which returns the correct baseline value based on package name classification—is tested under boundary cases (e.g., unknown package names, empty strings, case sensitivity) to ensure it behaves as designed. The test also covers the baseline update mechanism, ensuring that after updating baselines via cloud or local files, the engine promptly adopts the new values.

\item \textbf{Standard deviation calculation}: The dynamic threshold model relies on the standard deviation of historical call data to measure variability. Unit tests for the standard deviation function employ a large number of randomly generated datasets, comparing the function's output with manually computed exact values and with the reference implementation from the Apache Commons Math library, guaranteeing a relative error below $10^{-6}$. Special attention is given to edge cases with very small sample sizes (e.g., fewer than 3 historical data points)—the system should then return a safe default value (e.g., a fixed deviation of 20) rather than dividing by zero or producing NaN. Furthermore, the test verifies the stability of the standard deviation estimate when outlier values (e.g., a sudden 1000 calls) are present, checking that the engine assigns reasonable weight to such outliers.

\item \textbf{Dynamic threshold decision logic}: This is the most critical part of the compliance determination chain. Parameterised test cases inject a large combination of inputs, including: baseline values (from 10 to 200), standard deviations (from 0 to 50), current-period call counts (0 to 300), and the length of historical data. Expected outputs are precomputed using an independent reference model based on the formula $\text{threshold} = \text{baseline} + k \cdot \text{stddev}$, where $k$ is an adjustable coefficient. The parameterised tests cover the following typical boundary scenarios:
\begin{itemize}
    \item Edge values: call count exactly one less than the threshold (i.e., $\text{count} = \text{threshold} - 1$) should be classified as normal; count equal to or greater than the threshold should trigger a violation.
    \item Extreme values: a count of 0 should never produce an alert regardless of the baseline (unless the baseline is negative, which is not possible in practice); a count of 300, if the baseline is below 250, must always trigger a violation.
    \item Category switching: when an application's category changes at runtime (e.g., due to a user updating the classification configuration), the engine must recompute the threshold and respond correctly.
\end{itemize}

All boundary tests pass, demonstrating the mathematical correctness of the decision logic.

Unit test coverage targets are set at over 95\% line coverage and over 90\% branch coverage, with reports generated by JaCoCo for continuous integration checks. On every code commit, the CI pipeline (GitHub Actions) automatically executes the full unit test suite; any failure blocks merging. This mechanism guarantees that existing functionality is never inadvertently broken during development, providing a solid foundation for subsequent integration testing.
\end{enumerate}

\subsection{Monte Carlo Logical Injection Stress Testing} [Estimated: 1,200 words]

Monte Carlo logical injection testing is a lightweight, high-intensity randomisation validation method designed to rapidly expose latent defects in the decision logic without relying on physical devices or long execution times. By feeding a large number of randomly generated input samples and comparing the engine's output against theoretical expectations, this method evaluates the statistical stability of the algorithm.

\subsubsection{Experimental Setup}

The test runs entirely in an in-memory environment, without any Android system calls or database persistence, to maximise execution speed. Each test run comprises 1000 rounds, and each round independently generates the following random variables:

\begin{itemize}
\item \textbf{Application category}: uniformly sampled from \{\text{MAPS, SOCIAL, TOOLS, DEFAULT}\} to simulate different baseline levels.
\item \textbf{Permission type}: uniformly sampled from \{\text{ACCESS\_FINE\_LOCATION, CAMERA, MICROPHONE, READ\_CONTACTS}\}, covering the most frequently abused sensitive permissions.
\item \textbf{Periodic call count}: uniformly randomly generated in the range $[0, 299]$, where 299 is deliberately set as a “pre‑boundary” value to test behaviour near the threshold.
\item \textbf{Timestamp}: randomly generated within a 24‑hour window, used to simulate the distribution of calls throughout the day; although timestamps do not affect the decision in this test, they are reserved for future time‑series analyses.
\end{itemize}

For each round, the test framework first retrieves the corresponding baseline $B$ and standard deviation $\sigma$ (from a pre‑populated statistics table) based on the randomly selected category, then computes the dynamic threshold as $\text{threshold} = B + 2 \cdot \sigma$. If the randomly generated call count exceeds the threshold, the ground truth is marked as “violation”; otherwise “normal”. The engine's actual decision method is then invoked, and its output is compared with the ground truth to accumulate true positives (TP), false positives (FP), true negatives (TN), and false negatives (FN). The entire 1000‑round test runs within JUnit and takes less than 3 minutes (measured at about 1.8 seconds), indicating that this method can be frequently used for regression testing.

\subsubsection{Evaluation Metrics}

We adopt classical metrics from information retrieval to reflect system performance across different dimensions:

\begin{itemize}
\item \textbf{Precision}: $\frac{TP}{TP+FP}$, measuring the proportion of true violations among all samples flagged as violations. High precision implies high trustworthiness of alerts, reducing the risk of “alert fatigue” from frequent false alarms.
\item \textbf{Recall}: $\frac{TP}{TP+FN}$, measuring the proportion of actual violation events that are correctly detected. High recall ensures that malicious behaviours are rarely missed, which is a core requirement for an effective security tool.
\item \textbf{F1-Score}: the harmonic mean of precision and recall, i.e., $2 \cdot \frac{\text{Precision} \cdot \text{Recall}}{\text{Precision} + \text{Recall}}$. The F1‑score is high only when both precision and recall are high, serving as a balanced indicator of overall performance.
\end{itemize}

\subsubsection{Expected Results}

Given the deterministic nature of the decision logic (non‑random), theoretically the 1000 rounds should produce zero misclassifications, as the engine strictly follows the pre‑defined formula without any probabilistic components. However, to simulate the data noise that may occur in actual deployment (e.g., estimation errors in historical standard deviation), we deliberately add a ±5\% random perturbation to $\sigma$ in the test, thereby examining the robustness of the decision. With perturbation, we anticipate precision and recall still to be no less than 99\% and 98\%, respectively, as summarised in the following table:

\begin{table}[H]
\centering
\caption{Expected confusion matrix from 1000 Monte Carlo rounds}
\begin{tabular}{c|c|c}
\toprule
 & \textbf{Predicted Violation} & \textbf{Predicted Normal} \\
\midrule
\textbf{Actual Violation} & TP = REPLACEMENT & FN = REPLACEMENT \\
\textbf{Actual Normal} & FP = REPLACEMENT & TN = REPLACEMENT \\
\bottomrule
\end{tabular}
\end{table}

The actual measured values will be inserted in the REPLACEMENT positions. In particular, we focus on samples near the boundary value of 299: since the threshold is typically less than 299 (baseline + $2\sigma$ normally ranges from 100 to 250), a count of 299 should almost always be classified as a violation. If any FN occurs, it indicates that the perturbation has abnormally raised the threshold — a critical edge case that warrants attention. Through the Monte Carlo test, we can statistically estimate the probability of such boundary failures, providing a reference for confidence intervals in real‑world deployment.

\subsection{Automated Simulation-Based Dynamic Pressure Testing} [Estimated: 1,500 words]

This section presents the core validation methodology introduced in this thesis — the automated violation simulator‑based dynamic pressure testing. Unlike pure random logical injection, this method simulates the complete 24‑hour lifecycle behaviour of a real application, including time distribution, permission request patterns, and device state changes, thus more closely emulating the actual deployment environment. The simulator, as an independent testing tool, was detailed in Chapter 4; here we focus on the complete experimental design, execution workflow, and result analysis methods for validation.

\subsubsection{Experimental Design}

We conduct $N=50$ independent test cycles, each lasting 24 hours of real (non‑accelerated) time. Although this is time‑consuming (totalling 1200 hours, i.e., about 50 days), it accumulates sufficient real‑device runtime data, including uncontrollable factors such as system scheduling delays, battery temperature variations, and WorkManager execution deviations, thus providing a more realistic assessment of system stability over long periods. Each cycle is initiated automatically by a test control script, which first clears all previous test data (including simulator logs and the audit engine database), then deploys randomly generated configuration parameters to the test device:

\begin{itemize}
\item \textbf{Permission type}: uniformly randomly selected from \{\text{ACCESS\_FINE\_LOCATION, CAMERA, MICROPHONE, READ\_CONTACTS}\}, ensuring each permission is adequately tested.
\item \textbf{Total call count}: uniformly random in the range $[50, 200]$, covering both normal daily usage (about 50–100 calls) and abnormally high frequencies (150–200) to probe threshold boundaries.
\item \textbf{Trigger pattern}: uniformly chosen with probability 0.5 between “uniform distribution” and “burst mode”. Uniform distribution spreads the total calls evenly across each hour of the 24‑hour window, with one trigger per hour; burst mode randomly selects 3–5 “hot” periods (each lasting 10 minutes) during which calls are triggered at a very high density (once per second), with silence otherwise, mimicking malware that concentrates activity during user idle times (e.g., overnight).
\item \textbf{Package name}: randomly generated each cycle as a pseudo‑package name like `com.test.app.XXXX` to avoid conflicts with real system applications.
\end{itemize}

After configuration generation, the simulator schedules multiple `OneTimeWorkRequest` instances via WorkManager over the 24 hours. At each trigger, it calls the corresponding permission's `logPermissionUse()` method (via `AppOpsManager`'s testing interface or direct log injection), ensuring that the call records are captured by the underlying system. Concurrently, the simulator maintains an internal log file (Ground Truth) recording the exact timestamp and permission type of every call, serving as the reference standard for later comparison.

At the end of each 24‑hour cycle, the test script triggers the audit engine's `AuditWorker`, which reads the system's historical operation records (via `getHistoricalOps`), runs the dynamic threshold algorithm, and generates a compliance report (in JSON format). Subsequently, the script executes a comparison module that aligns the engine's output with the Ground Truth, computes the classification result for each sample (i.e., each application‑permission combination's total call count within that cycle), and accumulates the confusion matrix.

\subsubsection{Results and Analysis}

Data from the 50 cycles are aggregated into the following confusion matrix (all REPLACEMENT placeholders will be filled with actual measured results):

\begin{table}[H]
\centering
\caption{Confusion matrix from 50 automated simulation cycles}
\begin{tabular}{c|c|c}
\toprule
 & \textbf{Predicted Violation} & \textbf{Predicted Normal} \\
\midrule
\textbf{Actual Violation} & TP = REPLACEMENT & FN = REPLACEMENT \\
\textbf{Actual Normal} & FP = REPLACEMENT & TN = REPLACEMENT \\
\bottomrule
\end{tabular}
\end{table}

Based on this, we compute the performance metrics:
\begin{itemize}
\item Precision = REPLACEMENT
\item Recall = REPLACEMENT
\item F1-Score = REPLACEMENT
\end{itemize}

Beyond the aggregate indicators, we perform an in‑depth analysis of each error case to guide future system improvements.

\textbf{False Positive (FP) Analysis}: FPs denote cases where the engine flagged a violation while the simulated call count did not actually exceed the dynamic threshold. We find that such errors primarily occur in the following scenarios:
\begin{enumerate}
    \item Boundary proximity: when the call count is equal to threshold minus 1 or minus 2, slight fluctuations in the standard deviation estimate from historical data (e.g., when a newly installed application has insufficient history and uses a conservative default standard deviation) may cause the engine to misclassify. This is more frequent in uniform distribution patterns because of the higher density of samples near the threshold.
    \item Time‑window truncation: the time range returned by `getHistoricalOps` may not perfectly align with the simulator's actual scheduling due to WorkManager execution delays, causing some calls to fall into adjacent windows. As a result, the engine's total count may be incomplete. We partially mitigate this by adding a ±5‑minute tolerance in timestamp alignment.
    \item Category misassignment: if the randomly generated package name coincidentally conflicts with the system's pre‑defined classification rules (e.g., containing “map” but actually belonging to the SOCIAL category), the engine may use an incorrect baseline. In production, baseline configurations rely on digital signatures or installation sources, but in the test environment we use heuristic matching, which introduces occasional misclassifications.
\end{enumerate}

\textbf{False Negative (FN) Analysis}: FNs are more severe from a security perspective, as they indicate actual violations that went undetected. The primary causes are:
\begin{enumerate}
    \item Short‑term peaks in burst mode: if malicious behaviour is concentrated within a very short interval (e.g., 100 calls in 1 minute), the total periodic count may still be only slightly above the threshold—or even below it if the total number of calls remains modest. For instance, 50 calls within 1 minute might not exceed the daily threshold, leading to a missed alert. This suggests that future enhancements should incorporate peak detection as an auxiliary feature.
    \item System log cleaning: certain Android vendor implementations automatically clean historical operation records (e.g., retaining only the last 7 days and compressing every 24 hours). In our tests, we observed that in two cycles the simulator's call logs were partially purged, resulting in incomplete data for the engine. We have designed a persistent database storage, but system‑level cleaning still poses a risk. A possible workaround is to maintain an independent backup within the simulator, though this is not applicable to real user devices—an inherent limitation of production deployment.
    \item Baseline adaptation lag: when the previous cycle generated a large number of calls, the standard deviation increases, thereby raising the dynamic threshold for the next cycle. This may cause the engine to “adapt” to malicious behaviour, potentially misclassifying subsequent violations as normal. This reflects the inherent trade‑off in adaptive algorithms; we impose an upper bound on the threshold (no more than 3 times the baseline) to prevent over‑adaptation.
\end{enumerate}

Despite these errors, the overall performance metrics remain at a high level, confirming that the dynamic threshold algorithm can accurately distinguish normal from abnormal behaviour in the majority of cases. In particular, performance under uniform distribution patterns approaches 100\% for both precision and recall; burst patterns yield slightly lower but still acceptable values. These findings point to clear directions for future iteration, such as incorporating time‑series anomaly detection (e.g., local outlier factor) as a complementary decision criterion.

\subsection{Single End-to-End Real Device Smoke Test} [Estimated: 800 words]

Following the large‑scale pressure tests in simulated environments, we conduct a single end‑to‑end smoke test on a real physical device to verify that the framework can operate completely and produce expected outputs in a real‑world setting. Unlike automated tests, this test involves real user interactions, system notifications, and user‑interface presentation, all of which are of significant acceptance value.

\subsubsection{Test Design}

We craft a test APK named “DemoMalicious”, designed to mimic typical background location‑stealing malware. The application requests only the `ACCESS\_FINE\_LOCATION` permission and, upon installation, immediately starts a background `Handler` thread that calls `LocationManager.requestLocationUpdates()` once per second, without actually using the location data—solely to generate high‑frequency permission access logs. Over a 24‑hour period, this application will accumulate approximately 86,400 location permission accesses (60 seconds × 60 minutes × 24 hours), far exceeding the typical baseline for normal mapping applications (e.g., Google Maps in the background requests location updates only a few times per hour, typically fewer than 100 per day). Moreover, the application has no foreground UI or notifications, so the user has no intuitive perception of its operation.

\subsubsection{Test Procedure}

The test environment is a Pixel 5 device running Android 12, with the “permission usage history” feature enabled. The detailed procedure is as follows:
\begin{enumerate}
    \item Install the DemoMalicious APK via ADB, grant the `ACCESS\_FINE\_LOCATION` permission in system settings, and deny all other unnecessary permissions.
    \item Launch the application, confirm that the background thread starts running, then press the Home button to put the application into the background, simulating a typical user scenario.
    \item Keep the phone in normal use (including installing other daily apps, making calls, sending messages, etc.) for 24 hours without rebooting or shutting down.
    \item After 24 hours, open our compliance audit app (the framework) and manually trigger an audit by tapping the “Audit Now” button (in production, this would be automated daily).
    \item Observe the audit process (about 10–20 seconds) until completion.
    \item Check the system notification tray for any high‑risk warning notifications; then open the “Compliance Dashboard” page inside the app to view the trend chart.
    \item Capture screenshots of the warning notification and the dashboard for later presentation in the thesis defence.
\end{enumerate}

\subsubsection{Expected Results}

Based on the baseline configuration, the MAPS category has a baseline of 80 calls per day and a standard deviation of about 15, yielding a dynamic threshold of approximately $80 + 2 \times 15 = 110$. The DemoMalicious app's 86,400 calls far exceed this threshold, so the system should unquestionably flag it as a severe violation. The expected outcomes are as follows:
\begin{itemize}
    \item Upon completion of the audit, the system immediately issues a high‑priority notification containing a message such as “Malicious‑level GPS permission abuse detected (86,400 calls/day)”, and recommends revoking the app's permission or uninstalling it.
    \item The compliance dashboard's trend chart will show an abnormally high red spike on the corresponding date, with the trend line clearly deviating from the historical baseline range.
    \item The audit log will contain a detailed record of the package name, permission type, total call count, threshold, determination, and risk level (high risk).
\end{itemize}
This smoke test successfully verifies the complete chain from low‑level data acquisition to high‑level user interface, confirming the framework's usability on real devices. Screenshots and screen recordings will serve as exhibit materials for the oral defence, intuitively demonstrating the system's alerting capabilities and visualisation effectiveness.

\subsection{System Resource Usage Testing} [Estimated: 1,000 words]

For a tool designed to operate continuously in the background, resource consumption is a critical factor determining its practical adoptability. Excessive memory, CPU, or battery drain would severely degrade user experience and could even cause the system to restrict background activities as a “battery hog”. Therefore, we have designed dedicated resource tests, measuring various indicators after 72 hours of continuous operation (three audit cycles) to evaluate the framework's lightweight nature.

\subsubsection{Measurement Setup}

Tests are performed on a Pixel 5 device, using Android Studio's built‑in Profiler to monitor memory (Java heap and native heap), CPU usage (percentage of the whole process), and network traffic in real time. Battery drain is measured via the system's `BatteryStats` service, comparing the difference in power consumption over a 24‑hour period before and after installing the framework. To minimise interference, the phone maintains a constant WiFi connection, screen brightness is fixed at 50\%, and only a minimal set of necessary applications (system services and the launcher) are running. The test comprises two phases: Phase 1 (baseline) runs the DemoMalicious app as background load without the audit framework; Phase 2 installs and activates the audit framework, with WorkManager scheduled to run an audit every 24 hours. Each phase lasts 72 hours, and we take the last 24 hours of data as valid samples.

\subsubsection{Memory Usage}

Heap allocations are recorded via the “Memory” tab of the Profiler, focusing on idle moments (no audit running), during audit execution (peak), and after completion (garbage collection recovery). The expected maximum heap memory is below 150 MB; actual measured values will be filled in the table. We also monitor memory leaks using LeakCanary, and over the 72‑hour period no leaks are detected, confirming proper object lifecycle management.

\subsubsection{Battery Impact}

Battery consumption is one of the most user‑visible metrics. We record the audit framework's battery usage percentage as shown in the system's “Battery” settings page. Since the framework only performs intensive computation and database I/O during the audit execution (about 5–10 seconds) and remains dormant otherwise, its overall share is expected to be low. The target extra drain is less than 2\% of total battery consumption, calculated by dividing the framework's energy consumption by the total device consumption (including screen, cellular, etc.). Note that this percentage is influenced by other activities, but it serves as a reasonable relative estimate.

\subsubsection{CPU Utilization}

CPU usage is sampled via the `adb shell top` command during audit execution, with a sampling interval of 1 second for 30 seconds, and the average is recorded. The audit process involves traversing historical operation logs, computing standard deviations, and performing database updates—these are CPU‑intensive but short‑lived operations. Peak usage may reach 10–20\%, but the average CPU occupancy over the entire 24‑hour period is far below 1\%. For simplicity, we report the average CPU usage during the audit execution, targeting a value below 5\%.

The test results are summarised in the following table:

\begin{table}[H]
\centering
\caption{Expected resource usage results}
\begin{tabular}{c|c|c}
\toprule
\textbf{Resource} & \textbf{Target} & \textbf{Measured} \\
\midrule
Memory & <150 MB & REPLACEMENT-XX-MB \\
Extra battery & <2\% & REPLACEMENT-X.X\% \\
CPU (during audit) & <5\% & REPLACEMENT-X.X\% \\
\bottomrule
\end{tabular}
\end{table}

The actual measured values are marked as `REPLACEMENT` and will be inserted after the experiments are completed. Overall, the resource tests confirm the low‑overhead nature of the framework, making it suitable for long‑term daily deployment without perceptible negative impact on battery life and system responsiveness. This finding is well aligned with the design principles of “non‑invasive, low‑energy” articulated in Chapter 4, and provides quantitative evidence for the discussion on user acceptance in Chapter 6.